% --- Archon physics formalization source begin ---
% archon:physics
% archon:covers PhyXMiniProblems/problem_phyx_mini_0912.lean
% archon:source-report reports/phyx_mini/problem_phyx_mini_0912.source.json

\chapter{Physics problem phyx\_mini\_0912}
\label{ch:PhyXMiniProblems_problem_phyx_mini_0912}

\paragraph{Problem source.}
\#\# Physical scenario

In a vacuum chamber, a proton orbits a \textbackslash{}(1.0\textbackslash{},\textbackslash{}mathrm\{cm\}\textbackslash{})-diameter metal ball \textbackslash{}(1.0\textbackslash{},\textbackslash{}mathrm\{mm\}\textbackslash{}) above the surface with a period of \textbackslash{}(1.0\textbackslash{},\textbackslash{}mu\textbackslash{}mathrm\{s\}\textbackslash{}).

\#\# Question

What is the charge on the ball?

\#\# Answer choices

- A: A: -1.3 \textbackslash{}times 10\textasciicircum{}\{-12\}C

- B: B: 5.3 \textbackslash{}times 10\textasciicircum{}\{-14\}C

- C: C: -9.9 \textbackslash{}times 10\textasciicircum{}\{-12\}C

- D: D: 3.5 \textbackslash{}times 10\textasciicircum{}\{-14\}C

\#\# Figure caption (auxiliary; use the image as primary evidence)

The image is a diagram representing a charged particle system. Here's a detailed description:

- **Central Charge:** A large central sphere is labeled with a positive charge "Q."
- **Electron:** Outside the sphere, there is a smaller circle labeled with a positive sign and "e." This represents an electron.
- **Electric Field Lines (E):** Several arrows point radially inward towards the central charge, labeled as \textbackslash{}( \textbackslash{}vec\{E\} \textbackslash{}), indicating the direction of the electric field.
- **Force on Electron (\textbackslash{}( \textbackslash{}vec\{F\}\_\{\textbackslash{}text\{elec\}\} \textbackslash{})):** An arrow labeled \textbackslash{}( \textbackslash{}vec\{F\}\_\{\textbackslash{}text\{elec\}\} \textbackslash{}) points towards the central charge from the electron, showing the direction of the electric force on the electron.
- **Velocity (\textbackslash{}( \textbackslash{}vec\{v\} \textbackslash{})):** An arrow labeled \textbackslash{}( \textbackslash{}vec\{v\} \textbackslash{}) points perpendicularly to \textbackslash{}( \textbackslash{}vec\{F\}\_\{\textbackslash{}text\{elec\}\} \textbackslash{}), indicating the velocity of the electron.
- **Distance (r):** A line extends from the center of the sphere to the electron, labeled "r," representing the distance between the central charge and the electron. 

The diagram illustrates the electrostatic interaction between the central charge and an electron in its electric field.

\#\# Dataset metadata

Electrostatics; Multi-Formula Reasoning, Physical Model Grounding Reasoning

\paragraph{Recorded answer/context.}
C: C: -9.9 \textbackslash{}times 10\textasciicircum{}\{-12\}C

\paragraph{Figure/image path.}
/root/proposal\_for\_physic/science-mango/phyx\_mini\_run/phyx\_data/test\_image/912.png

\paragraph{Formalization target.}
create a compiling Lean file with sorry bodies at `PhyXMiniProblems/problem\_phyx\_mini\_0912.lean`. The Lean declarations must preserve the physical quantities, dimensions or dimensional roles, figure labels, governing-law hypotheses, and final relation expressed by this problem.
Use Mathlib/Physlib names found through LeanExplore where available. If a physics API is missing, introduce faithful local abstractions rather than scalar placeholder aliases.

\begin{theorem}[Physics formalization target]
\label{thm:physics:phyx_mini_0912:target}
\lean{PhyXMiniProblems.ProblemPhyXMini0912.problem_phyx_mini_0912}
\uses{def:physics:phyx-mini-0912:phyxminiproblems-problemphyxmini0912-chargeincoulombs, def:physics:phyx-mini-0912:phyxminiproblems-problemphyxmini0912-protonorbitsetup, def:physics:phyx-mini-0912:phyxminiproblems-problemphyxmini0912-matcheswrittenscenario, def:physics:phyx-mini-0912:phyxminiproblems-problemphyxmini0912-matchesproblemreadouts, def:physics:phyx-mini-0912:phyxminiproblems-problemphyxmini0912-matchesprimaryorbitfigure, def:physics:phyx-mini-0912:phyxminiproblems-problemphyxmini0912-satisfiesorbitradiusgeometry, def:physics:phyx-mini-0912:phyxminiproblems-problemphyxmini0912-usestextbookprotonandcoulombconstants, def:physics:phyx-mini-0912:phyxminiproblems-problemphyxmini0912-hasphysicalprotonorbitparameters, def:physics:phyx-mini-0912:phyxminiproblems-problemphyxmini0912-satisfiesexteriorsphericalcoulombfieldlaw, def:physics:phyx-mini-0912:phyxminiproblems-problemphyxmini0912-satisfieselectricforcelawonproton, def:physics:phyx-mini-0912:phyxminiproblems-problemphyxmini0912-satisfiesuniformcircularorbitlaws, def:physics:phyx-mini-0912:phyxminiproblems-problemphyxmini0912-answerchoice, def:physics:phyx-mini-0912:phyxminiproblems-problemphyxmini0912-recordeddatasetanswer, def:physics:phyx-mini-0912:phyxminiproblems-problemphyxmini0912-ballchargeroundstochoicec, def:physics:phyx-mini-0912:phyxminiproblems-problemphyxmini0912-isclosestdisplayedcharge}
The assigned autoformalize agent should translate this physics problem into Lean declarations in the covered file, with theorem and lemma proof bodies written as `by sorry`.
\end{theorem}
\begin{proof}
This is an autoformalization task, not a proof task. Produce statements that can later be proved by the physics prover without weakening the physical model.
\end{proof}

% --- Archon declaration topology begin ---
\section{Declaration topology}
\begin{definition}[force Dimension]
  \label{def:physics:phyx-mini-0912:phyxminiproblems-problemphyxmini0912-forcedimension}
  \lean{PhyXMiniProblems.ProblemPhyXMini0912.forceDimension}
  The physical dimension M L T⁻² of force
\end{definition}
\begin{proof}
  This is the declared model definition of force Dimension.
\end{proof}
\begin{definition}[electric Field Strength Dimension]
  \label{def:physics:phyx-mini-0912:phyxminiproblems-problemphyxmini0912-electricfieldstrengthdimension}
  \lean{PhyXMiniProblems.ProblemPhyXMini0912.electricFieldStrengthDimension}
  \uses{def:physics:phyx-mini-0912:phyxminiproblems-problemphyxmini0912-forcedimension}
  The physical dimension M L T⁻² C⁻¹ of electric-field strength
\end{definition}
\begin{proof}
  This is the declared model definition of electric Field Strength Dimension.
\end{proof}
\begin{definition}[velocity Dimension]
  \label{def:physics:phyx-mini-0912:phyxminiproblems-problemphyxmini0912-velocitydimension}
  \lean{PhyXMiniProblems.ProblemPhyXMini0912.velocityDimension}
  The physical dimension L T⁻¹ of velocity
\end{definition}
\begin{proof}
  This is the declared model definition of velocity Dimension.
\end{proof}
\begin{definition}[Length Quantity]
  \label{def:physics:phyx-mini-0912:phyxminiproblems-problemphyxmini0912-lengthquantity}
  \lean{PhyXMiniProblems.ProblemPhyXMini0912.LengthQuantity}
  A nonnegative, unit-independent physical length
\end{definition}
\begin{proof}
  This is the declared model definition of Length Quantity.
\end{proof}
\begin{definition}[Duration Quantity]
  \label{def:physics:phyx-mini-0912:phyxminiproblems-problemphyxmini0912-durationquantity}
  \lean{PhyXMiniProblems.ProblemPhyXMini0912.DurationQuantity}
  A nonnegative, unit-independent duration
\end{definition}
\begin{proof}
  This is the declared model definition of Duration Quantity.
\end{proof}
\begin{definition}[Mass Quantity]
  \label{def:physics:phyx-mini-0912:phyxminiproblems-problemphyxmini0912-massquantity}
  \lean{PhyXMiniProblems.ProblemPhyXMini0912.MassQuantity}
  A nonnegative, unit-independent mass
\end{definition}
\begin{proof}
  This is the declared model definition of Mass Quantity.
\end{proof}
\begin{definition}[Signed Charge Quantity]
  \label{def:physics:phyx-mini-0912:phyxminiproblems-problemphyxmini0912-signedchargequantity}
  \lean{PhyXMiniProblems.ProblemPhyXMini0912.SignedChargeQuantity}
  A signed, unit-independent electric charge
\end{definition}
\begin{proof}
  This is the declared model definition of Signed Charge Quantity.
\end{proof}
\begin{definition}[Electric Field Magnitude Quantity]
  \label{def:physics:phyx-mini-0912:phyxminiproblems-problemphyxmini0912-electricfieldmagnitudequantity}
  \lean{PhyXMiniProblems.ProblemPhyXMini0912.ElectricFieldMagnitudeQuantity}
  \uses{def:physics:phyx-mini-0912:phyxminiproblems-problemphyxmini0912-electricfieldstrengthdimension}
  A nonnegative, unit-independent electric-field magnitude
\end{definition}
\begin{proof}
  This is the declared model definition of Electric Field Magnitude Quantity.
\end{proof}
\begin{definition}[Spatial Velocity Quantity]
  \label{def:physics:phyx-mini-0912:phyxminiproblems-problemphyxmini0912-spatialvelocityquantity}
  \lean{PhyXMiniProblems.ProblemPhyXMini0912.SpatialVelocityQuantity}
  \uses{def:physics:phyx-mini-0912:phyxminiproblems-problemphyxmini0912-velocitydimension}
  A unit-independent spatial velocity vector
\end{definition}
\begin{proof}
  This is the declared model definition of Spatial Velocity Quantity.
\end{proof}
\begin{definition}[Spatial Force Quantity]
  \label{def:physics:phyx-mini-0912:phyxminiproblems-problemphyxmini0912-spatialforcequantity}
  \lean{PhyXMiniProblems.ProblemPhyXMini0912.SpatialForceQuantity}
  \uses{def:physics:phyx-mini-0912:phyxminiproblems-problemphyxmini0912-forcedimension}
  A unit-independent spatial force vector
\end{definition}
\begin{proof}
  This is the declared model definition of Spatial Force Quantity.
\end{proof}
\begin{definition}[length In Meters]
  \label{def:physics:phyx-mini-0912:phyxminiproblems-problemphyxmini0912-lengthinmeters}
  \lean{PhyXMiniProblems.ProblemPhyXMini0912.lengthInMeters}
  \uses{def:physics:phyx-mini-0912:phyxminiproblems-problemphyxmini0912-lengthquantity}
  Coherent-SI readout of a length, in metres
\end{definition}
\begin{proof}
  This is the declared model definition of length In Meters.
\end{proof}
\begin{definition}[length In Centimeters]
  \label{def:physics:phyx-mini-0912:phyxminiproblems-problemphyxmini0912-lengthincentimeters}
  \lean{PhyXMiniProblems.ProblemPhyXMini0912.lengthInCentimeters}
  \uses{def:physics:phyx-mini-0912:phyxminiproblems-problemphyxmini0912-lengthquantity, def:physics:phyx-mini-0912:phyxminiproblems-problemphyxmini0912-lengthinmeters}
  Centimetre readout used for the ball-diameter datum
\end{definition}
\begin{proof}
  This is the declared model definition of length In Centimeters.
\end{proof}
\begin{definition}[length In Millimeters]
  \label{def:physics:phyx-mini-0912:phyxminiproblems-problemphyxmini0912-lengthinmillimeters}
  \lean{PhyXMiniProblems.ProblemPhyXMini0912.lengthInMillimeters}
  \uses{def:physics:phyx-mini-0912:phyxminiproblems-problemphyxmini0912-lengthquantity, def:physics:phyx-mini-0912:phyxminiproblems-problemphyxmini0912-lengthinmeters}
  Millimetre readout used for the proton's surface clearance
\end{definition}
\begin{proof}
  This is the declared model definition of length In Millimeters.
\end{proof}
\begin{definition}[duration In Seconds]
  \label{def:physics:phyx-mini-0912:phyxminiproblems-problemphyxmini0912-durationinseconds}
  \lean{PhyXMiniProblems.ProblemPhyXMini0912.durationInSeconds}
  \uses{def:physics:phyx-mini-0912:phyxminiproblems-problemphyxmini0912-durationquantity}
  Coherent-SI readout of a duration, in seconds
\end{definition}
\begin{proof}
  This is the declared model definition of duration In Seconds.
\end{proof}
\begin{definition}[duration In Microseconds]
  \label{def:physics:phyx-mini-0912:phyxminiproblems-problemphyxmini0912-durationinmicroseconds}
  \lean{PhyXMiniProblems.ProblemPhyXMini0912.durationInMicroseconds}
  \uses{def:physics:phyx-mini-0912:phyxminiproblems-problemphyxmini0912-durationquantity, def:physics:phyx-mini-0912:phyxminiproblems-problemphyxmini0912-durationinseconds}
  Microsecond readout used for the orbital-period datum
\end{definition}
\begin{proof}
  This is the declared model definition of duration In Microseconds.
\end{proof}
\begin{definition}[mass In Kilograms]
  \label{def:physics:phyx-mini-0912:phyxminiproblems-problemphyxmini0912-massinkilograms}
  \lean{PhyXMiniProblems.ProblemPhyXMini0912.massInKilograms}
  \uses{def:physics:phyx-mini-0912:phyxminiproblems-problemphyxmini0912-massquantity}
  Coherent-SI readout of a mass, in kilograms
\end{definition}
\begin{proof}
  This is the declared model definition of mass In Kilograms.
\end{proof}
\begin{definition}[charge In Coulombs]
  \label{def:physics:phyx-mini-0912:phyxminiproblems-problemphyxmini0912-chargeincoulombs}
  \lean{PhyXMiniProblems.ProblemPhyXMini0912.chargeInCoulombs}
  \uses{def:physics:phyx-mini-0912:phyxminiproblems-problemphyxmini0912-signedchargequantity}
  Coherent-SI readout of a signed charge, in coulombs
\end{definition}
\begin{proof}
  This is the declared model definition of charge In Coulombs.
\end{proof}
\begin{definition}[electric Field Magnitude In Newtons Per Coulomb]
  \label{def:physics:phyx-mini-0912:phyxminiproblems-problemphyxmini0912-electricfieldmagnitudeinnewtonspercoulomb}
  \lean{PhyXMiniProblems.ProblemPhyXMini0912.electricFieldMagnitudeInNewtonsPerCoulomb}
  \uses{def:physics:phyx-mini-0912:phyxminiproblems-problemphyxmini0912-electricfieldmagnitudequantity}
  Coherent-SI readout of electric-field magnitude, in newtons per coulomb
\end{definition}
\begin{proof}
  This is the declared model definition of electric Field Magnitude In Newtons Per Coulomb.
\end{proof}
\begin{definition}[velocity Vector In Meters Per Second]
  \label{def:physics:phyx-mini-0912:phyxminiproblems-problemphyxmini0912-velocityvectorinmeterspersecond}
  \lean{PhyXMiniProblems.ProblemPhyXMini0912.velocityVectorInMetersPerSecond}
  \uses{def:physics:phyx-mini-0912:phyxminiproblems-problemphyxmini0912-spatialvelocityquantity}
  Coherent-SI readout of velocity, in metres per second
\end{definition}
\begin{proof}
  This is the declared model definition of velocity Vector In Meters Per Second.
\end{proof}
\begin{definition}[speed In Meters Per Second]
  \label{def:physics:phyx-mini-0912:phyxminiproblems-problemphyxmini0912-speedinmeterspersecond}
  \lean{PhyXMiniProblems.ProblemPhyXMini0912.speedInMetersPerSecond}
  \uses{def:physics:phyx-mini-0912:phyxminiproblems-problemphyxmini0912-spatialvelocityquantity, def:physics:phyx-mini-0912:phyxminiproblems-problemphyxmini0912-velocityvectorinmeterspersecond}
  Speed obtained as the norm of the coherent-SI velocity vector
\end{definition}
\begin{proof}
  This is the declared model definition of speed In Meters Per Second.
\end{proof}
\begin{definition}[force Vector In Newtons]
  \label{def:physics:phyx-mini-0912:phyxminiproblems-problemphyxmini0912-forcevectorinnewtons}
  \lean{PhyXMiniProblems.ProblemPhyXMini0912.forceVectorInNewtons}
  \uses{def:physics:phyx-mini-0912:phyxminiproblems-problemphyxmini0912-spatialforcequantity}
  Coherent-SI readout of force, in newtons
\end{definition}
\begin{proof}
  This is the declared model definition of force Vector In Newtons.
\end{proof}
\begin{definition}[force Magnitude In Newtons]
  \label{def:physics:phyx-mini-0912:phyxminiproblems-problemphyxmini0912-forcemagnitudeinnewtons}
  \lean{PhyXMiniProblems.ProblemPhyXMini0912.forceMagnitudeInNewtons}
  \uses{def:physics:phyx-mini-0912:phyxminiproblems-problemphyxmini0912-spatialforcequantity, def:physics:phyx-mini-0912:phyxminiproblems-problemphyxmini0912-forcevectorinnewtons}
  Force magnitude obtained as the norm of the coherent-SI force vector
\end{definition}
\begin{proof}
  This is the declared model definition of force Magnitude In Newtons.
\end{proof}
\begin{definition}[Chamber Medium]
  \label{def:physics:phyx-mini-0912:phyxminiproblems-problemphyxmini0912-chambermedium}
  \lean{PhyXMiniProblems.ProblemPhyXMini0912.ChamberMedium}
  Medium in which the charged-particle motion takes place
\end{definition}
\begin{proof}
  This is the declared model definition of Chamber Medium.
\end{proof}
\begin{definition}[Particle Species]
  \label{def:physics:phyx-mini-0912:phyxminiproblems-problemphyxmini0912-particlespecies}
  \lean{PhyXMiniProblems.ProblemPhyXMini0912.ParticleSpecies}
  Particle species relevant to the prose/figure discrepancy
\end{definition}
\begin{proof}
  This is the declared model definition of Particle Species.
\end{proof}
\begin{definition}[Central Body Model]
  \label{def:physics:phyx-mini-0912:phyxminiproblems-problemphyxmini0912-centralbodymodel}
  \lean{PhyXMiniProblems.ProblemPhyXMini0912.CentralBodyModel}
  Model assigned to the central charged body
\end{definition}
\begin{proof}
  This is the declared model definition of Central Body Model.
\end{proof}
\begin{definition}[Figure Arrow]
  \label{def:physics:phyx-mini-0912:phyxminiproblems-problemphyxmini0912-figurearrow}
  \lean{PhyXMiniProblems.ProblemPhyXMini0912.FigureArrow}
  Named vector arrows visible in image 912.png
\end{definition}
\begin{proof}
  This is the declared model definition of Figure Arrow.
\end{proof}
\begin{definition}[Vector Direction]
  \label{def:physics:phyx-mini-0912:phyxminiproblems-problemphyxmini0912-vectordirection}
  \lean{PhyXMiniProblems.ProblemPhyXMini0912.VectorDirection}
  Qualitative direction of a physical or depicted vector
\end{definition}
\begin{proof}
  This is the declared model definition of Vector Direction.
\end{proof}
\begin{definition}[Charged Ball Orbit Figure]
  \label{def:physics:phyx-mini-0912:phyxminiproblems-problemphyxmini0912-chargedballorbitfigure}
  \lean{PhyXMiniProblems.ProblemPhyXMini0912.ChargedBallOrbitFigure}
  \uses{def:physics:phyx-mini-0912:phyxminiproblems-problemphyxmini0912-figurearrow, def:physics:phyx-mini-0912:phyxminiproblems-problemphyxmini0912-vectordirection}
  Literal content transcribed from the primary raster. The raster's orbiting marker is labelled e and contains a plus sign; it does not contain the word "electron". Charge values and answer choices do not occur in this structure
\end{definition}
\begin{proof}
  This is the declared model definition of Charged Ball Orbit Figure.
\end{proof}
\begin{definition}[Proton Orbit Setup]
  \label{def:physics:phyx-mini-0912:phyxminiproblems-problemphyxmini0912-protonorbitsetup}
  \lean{PhyXMiniProblems.ProblemPhyXMini0912.ProtonOrbitSetup}
  \uses{def:physics:phyx-mini-0912:phyxminiproblems-problemphyxmini0912-lengthquantity, def:physics:phyx-mini-0912:phyxminiproblems-problemphyxmini0912-durationquantity, def:physics:phyx-mini-0912:phyxminiproblems-problemphyxmini0912-massquantity, def:physics:phyx-mini-0912:phyxminiproblems-problemphyxmini0912-signedchargequantity, def:physics:phyx-mini-0912:phyxminiproblems-problemphyxmini0912-electricfieldmagnitudequantity, def:physics:phyx-mini-0912:phyxminiproblems-problemphyxmini0912-spatialvelocityquantity, def:physics:phyx-mini-0912:phyxminiproblems-problemphyxmini0912-spatialforcequantity, def:physics:phyx-mini-0912:phyxminiproblems-problemphyxmini0912-chambermedium, def:physics:phyx-mini-0912:phyxminiproblems-problemphyxmini0912-particlespecies, def:physics:phyx-mini-0912:phyxminiproblems-problemphyxmini0912-centralbodymodel, def:physics:phyx-mini-0912:phyxminiproblems-problemphyxmini0912-vectordirection, def:physics:phyx-mini-0912:phyxminiproblems-problemphyxmini0912-chargedballorbitfigure}
  Independent physical quantities in the proton-orbit experiment. In particular, ballCharge is not defined from an answer choice or from the recorded value -9.9 * 10\textasciicircum{}-12 C
\end{definition}
\begin{proof}
  This is the declared model definition of Proton Orbit Setup.
\end{proof}
\begin{definition}[Matches Written Scenario]
  \label{def:physics:phyx-mini-0912:phyxminiproblems-problemphyxmini0912-matcheswrittenscenario}
  \lean{PhyXMiniProblems.ProblemPhyXMini0912.MatchesWrittenScenario}
  \uses{def:physics:phyx-mini-0912:phyxminiproblems-problemphyxmini0912-protonorbitsetup}
  The prose's object and environment identifications
\end{definition}
\begin{proof}
  This is the declared model definition of Matches Written Scenario.
\end{proof}
\begin{definition}[Matches Problem Readouts]
  \label{def:physics:phyx-mini-0912:phyxminiproblems-problemphyxmini0912-matchesproblemreadouts}
  \lean{PhyXMiniProblems.ProblemPhyXMini0912.MatchesProblemReadouts}
  \uses{def:physics:phyx-mini-0912:phyxminiproblems-problemphyxmini0912-lengthincentimeters, def:physics:phyx-mini-0912:phyxminiproblems-problemphyxmini0912-lengthinmillimeters, def:physics:phyx-mini-0912:phyxminiproblems-problemphyxmini0912-durationinmicroseconds, def:physics:phyx-mini-0912:phyxminiproblems-problemphyxmini0912-protonorbitsetup}
  The three numerical readouts printed in the prose. They contain no charge readout for the ball
\end{definition}
\begin{proof}
  This is the declared model definition of Matches Problem Readouts.
\end{proof}
\begin{definition}[Matches Primary Orbit Figure]
  \label{def:physics:phyx-mini-0912:phyxminiproblems-problemphyxmini0912-matchesprimaryorbitfigure}
  \lean{PhyXMiniProblems.ProblemPhyXMini0912.MatchesPrimaryOrbitFigure}
  \uses{def:physics:phyx-mini-0912:phyxminiproblems-problemphyxmini0912-protonorbitsetup}
  Primary-raster evidence. Inward field and force arrows and the ball's negative marks are recorded as observations, while the physical reason for their directions is supplied separately by the Coulomb and force laws
\end{definition}
\begin{proof}
  This is the declared model definition of Matches Primary Orbit Figure.
\end{proof}
\begin{definition}[Satisfies Orbit Radius Geometry]
  \label{def:physics:phyx-mini-0912:phyxminiproblems-problemphyxmini0912-satisfiesorbitradiusgeometry}
  \lean{PhyXMiniProblems.ProblemPhyXMini0912.SatisfiesOrbitRadiusGeometry}
  \uses{def:physics:phyx-mini-0912:phyxminiproblems-problemphyxmini0912-lengthinmeters, def:physics:phyx-mini-0912:phyxminiproblems-problemphyxmini0912-protonorbitsetup}
  The center-to-proton radius is the ball radius plus the stated clearance. This is the geometric interpretation of "1.0 mm above the surface", not a charge conclusion
\end{definition}
\begin{proof}
  This is the declared model definition of Satisfies Orbit Radius Geometry.
\end{proof}
\begin{definition}[Uses Textbook Proton And Coulomb Constants]
  \label{def:physics:phyx-mini-0912:phyxminiproblems-problemphyxmini0912-usestextbookprotonandcoulombconstants}
  \lean{PhyXMiniProblems.ProblemPhyXMini0912.UsesTextbookProtonAndCoulombConstants}
  \uses{def:physics:phyx-mini-0912:phyxminiproblems-problemphyxmini0912-massinkilograms, def:physics:phyx-mini-0912:phyxminiproblems-problemphyxmini0912-chargeincoulombs, def:physics:phyx-mini-0912:phyxminiproblems-problemphyxmini0912-protonorbitsetup}
  Rounded school values used for the numerical multiple-choice calculation. They are independent calibration data and contain no value for ballCharge
\end{definition}
\begin{proof}
  This is the declared model definition of Uses Textbook Proton And Coulomb Constants.
\end{proof}
\begin{definition}[Has Physical Proton Orbit Parameters]
  \label{def:physics:phyx-mini-0912:phyxminiproblems-problemphyxmini0912-hasphysicalprotonorbitparameters}
  \lean{PhyXMiniProblems.ProblemPhyXMini0912.HasPhysicalProtonOrbitParameters}
  \uses{def:physics:phyx-mini-0912:phyxminiproblems-problemphyxmini0912-lengthinmeters, def:physics:phyx-mini-0912:phyxminiproblems-problemphyxmini0912-durationinseconds, def:physics:phyx-mini-0912:phyxminiproblems-problemphyxmini0912-massinkilograms, def:physics:phyx-mini-0912:phyxminiproblems-problemphyxmini0912-chargeincoulombs, def:physics:phyx-mini-0912:phyxminiproblems-problemphyxmini0912-electricfieldmagnitudeinnewtonspercoulomb, def:physics:phyx-mini-0912:phyxminiproblems-problemphyxmini0912-forcemagnitudeinnewtons, def:physics:phyx-mini-0912:phyxminiproblems-problemphyxmini0912-protonorbitsetup}
  Positivity and nondegeneracy conditions for the physical experiment
\end{definition}
\begin{proof}
  This is the declared model definition of Has Physical Proton Orbit Parameters.
\end{proof}
\begin{definition}[Satisfies Exterior Spherical Coulomb Field Law]
  \label{def:physics:phyx-mini-0912:phyxminiproblems-problemphyxmini0912-satisfiesexteriorsphericalcoulombfieldlaw}
  \lean{PhyXMiniProblems.ProblemPhyXMini0912.SatisfiesExteriorSphericalCoulombFieldLaw}
  \uses{def:physics:phyx-mini-0912:phyxminiproblems-problemphyxmini0912-lengthinmeters, def:physics:phyx-mini-0912:phyxminiproblems-problemphyxmini0912-chargeincoulombs, def:physics:phyx-mini-0912:phyxminiproblems-problemphyxmini0912-electricfieldmagnitudeinnewtonspercoulomb, def:physics:phyx-mini-0912:phyxminiproblems-problemphyxmini0912-protonorbitsetup}
  Exterior spherical Coulomb law for the charged conducting ball. The field magnitude is k |Q| / r², and its inward/outward orientation records the sign of the source charge. This is a general law and does not assign a numerical value to Q
\end{definition}
\begin{proof}
  This is the declared model definition of Satisfies Exterior Spherical Coulomb Field Law.
\end{proof}
\begin{definition}[Satisfies Electric Force Law On Proton]
  \label{def:physics:phyx-mini-0912:phyxminiproblems-problemphyxmini0912-satisfieselectricforcelawonproton}
  \lean{PhyXMiniProblems.ProblemPhyXMini0912.SatisfiesElectricForceLawOnProton}
  \uses{def:physics:phyx-mini-0912:phyxminiproblems-problemphyxmini0912-chargeincoulombs, def:physics:phyx-mini-0912:phyxminiproblems-problemphyxmini0912-electricfieldmagnitudeinnewtonspercoulomb, def:physics:phyx-mini-0912:phyxminiproblems-problemphyxmini0912-forcevectorinnewtons, def:physics:phyx-mini-0912:phyxminiproblems-problemphyxmini0912-forcemagnitudeinnewtons, def:physics:phyx-mini-0912:phyxminiproblems-problemphyxmini0912-protonorbitsetup}
  For the positive proton, the electric force has magnitude q E and points in the same direction as the electric field. No orbital-charge answer occurs in this force law
\end{definition}
\begin{proof}
  This is the declared model definition of Satisfies Electric Force Law On Proton.
\end{proof}
\begin{definition}[Satisfies Uniform Circular Orbit Laws]
  \label{def:physics:phyx-mini-0912:phyxminiproblems-problemphyxmini0912-satisfiesuniformcircularorbitlaws}
  \lean{PhyXMiniProblems.ProblemPhyXMini0912.SatisfiesUniformCircularOrbitLaws}
  \uses{def:physics:phyx-mini-0912:phyxminiproblems-problemphyxmini0912-lengthinmeters, def:physics:phyx-mini-0912:phyxminiproblems-problemphyxmini0912-durationinseconds, def:physics:phyx-mini-0912:phyxminiproblems-problemphyxmini0912-massinkilograms, def:physics:phyx-mini-0912:phyxminiproblems-problemphyxmini0912-velocityvectorinmeterspersecond, def:physics:phyx-mini-0912:phyxminiproblems-problemphyxmini0912-speedinmeterspersecond, def:physics:phyx-mini-0912:phyxminiproblems-problemphyxmini0912-forcemagnitudeinnewtons, def:physics:phyx-mini-0912:phyxminiproblems-problemphyxmini0912-protonorbitsetup}
  Uniform circular motion: the speed is circumference divided by period, the inward electric force supplies m v² / r, and the velocity is tangent to the orbit and perpendicular to the radius
\end{definition}
\begin{proof}
  This is the declared model definition of Satisfies Uniform Circular Orbit Laws.
\end{proof}
\begin{lemma}[orbit radius in meters]
  \label{lem:physics:phyx-mini-0912:phyxminiproblems-problemphyxmini0912-orbit-radius-in-meters}
  \lean{PhyXMiniProblems.ProblemPhyXMini0912.orbit_radius_in_meters}
  \uses{def:physics:phyx-mini-0912:phyxminiproblems-problemphyxmini0912-lengthinmeters, def:physics:phyx-mini-0912:phyxminiproblems-problemphyxmini0912-protonorbitsetup, def:physics:phyx-mini-0912:phyxminiproblems-problemphyxmini0912-matchesproblemreadouts, def:physics:phyx-mini-0912:phyxminiproblems-problemphyxmini0912-satisfiesorbitradiusgeometry}
  The stated diameter and surface clearance give an orbital radius of 6 mm
\end{lemma}
\begin{proof}
  The conclusion follows from the governing relations and figure readouts named in the statement of orbit radius in meters.
\end{proof}
\begin{lemma}[signed ball charge from circular orbit]
  \label{lem:physics:phyx-mini-0912:phyxminiproblems-problemphyxmini0912-signed-ball-charge-from-circular-orbit}
  \lean{PhyXMiniProblems.ProblemPhyXMini0912.signed_ball_charge_from_circular_orbit}
  \uses{def:physics:phyx-mini-0912:phyxminiproblems-problemphyxmini0912-lengthinmeters, def:physics:phyx-mini-0912:phyxminiproblems-problemphyxmini0912-durationinseconds, def:physics:phyx-mini-0912:phyxminiproblems-problemphyxmini0912-massinkilograms, def:physics:phyx-mini-0912:phyxminiproblems-problemphyxmini0912-chargeincoulombs, def:physics:phyx-mini-0912:phyxminiproblems-problemphyxmini0912-protonorbitsetup, def:physics:phyx-mini-0912:phyxminiproblems-problemphyxmini0912-matchesprimaryorbitfigure, def:physics:phyx-mini-0912:phyxminiproblems-problemphyxmini0912-hasphysicalprotonorbitparameters, def:physics:phyx-mini-0912:phyxminiproblems-problemphyxmini0912-satisfiesexteriorsphericalcoulombfieldlaw, def:physics:phyx-mini-0912:phyxminiproblems-problemphyxmini0912-satisfieselectricforcelawonproton, def:physics:phyx-mini-0912:phyxminiproblems-problemphyxmini0912-satisfiesuniformcircularorbitlaws}
  Eliminating the field, force, and speed from the governing laws gives the signed source-charge formula. Its leading minus sign is derived from the inward field and the sign clause of Coulomb's law
\end{lemma}
\begin{proof}
  The conclusion follows from the governing relations and figure readouts named in the statement of signed ball charge from circular orbit.
\end{proof}
\begin{definition}[Answer Choice]
  \label{def:physics:phyx-mini-0912:phyxminiproblems-problemphyxmini0912-answerchoice}
  \lean{PhyXMiniProblems.ProblemPhyXMini0912.AnswerChoice}
  Labels attached to the four displayed charge choices
\end{definition}
\begin{proof}
  This is the declared model definition of Answer Choice.
\end{proof}
\begin{definition}[answer Charge In Coulombs]
  \label{def:physics:phyx-mini-0912:phyxminiproblems-problemphyxmini0912-answerchargeincoulombs}
  \lean{PhyXMiniProblems.ProblemPhyXMini0912.answerChargeInCoulombs}
  \uses{def:physics:phyx-mini-0912:phyxminiproblems-problemphyxmini0912-answerchoice}
  Signed charge in coulombs printed beside each answer label
\end{definition}
\begin{proof}
  This is the declared model definition of answer Charge In Coulombs.
\end{proof}
\begin{definition}[recorded Dataset Answer]
  \label{def:physics:phyx-mini-0912:phyxminiproblems-problemphyxmini0912-recordeddatasetanswer}
  \lean{PhyXMiniProblems.ProblemPhyXMini0912.recordedDatasetAnswer}
  \uses{def:physics:phyx-mini-0912:phyxminiproblems-problemphyxmini0912-answerchoice}
  The answer label recorded by the supplied dataset
\end{definition}
\begin{proof}
  This is the declared model definition of recorded Dataset Answer.
\end{proof}
\begin{definition}[choice CDisplay Tolerance In Coulombs]
  \label{def:physics:phyx-mini-0912:phyxminiproblems-problemphyxmini0912-choicecdisplaytoleranceincoulombs}
  \lean{PhyXMiniProblems.ProblemPhyXMini0912.choiceCDisplayToleranceInCoulombs}
  Half of the last displayed decimal place in choice C
\end{definition}
\begin{proof}
  This is the declared model definition of choice CDisplay Tolerance In Coulombs.
\end{proof}
\begin{definition}[Ball Charge Rounds To Choice C]
  \label{def:physics:phyx-mini-0912:phyxminiproblems-problemphyxmini0912-ballchargeroundstochoicec}
  \lean{PhyXMiniProblems.ProblemPhyXMini0912.BallChargeRoundsToChoiceC}
  \uses{def:physics:phyx-mini-0912:phyxminiproblems-problemphyxmini0912-chargeincoulombs, def:physics:phyx-mini-0912:phyxminiproblems-problemphyxmini0912-protonorbitsetup, def:physics:phyx-mini-0912:phyxminiproblems-problemphyxmini0912-answerchargeincoulombs, def:physics:phyx-mini-0912:phyxminiproblems-problemphyxmini0912-choicecdisplaytoleranceincoulombs}
  The physical charge agrees with choice C to its displayed precision
\end{definition}
\begin{proof}
  This is the declared model definition of Ball Charge Rounds To Choice C.
\end{proof}
\begin{definition}[Is Closest Displayed Charge]
  \label{def:physics:phyx-mini-0912:phyxminiproblems-problemphyxmini0912-isclosestdisplayedcharge}
  \lean{PhyXMiniProblems.ProblemPhyXMini0912.IsClosestDisplayedCharge}
  \uses{def:physics:phyx-mini-0912:phyxminiproblems-problemphyxmini0912-chargeincoulombs, def:physics:phyx-mini-0912:phyxminiproblems-problemphyxmini0912-protonorbitsetup, def:physics:phyx-mini-0912:phyxminiproblems-problemphyxmini0912-answerchoice, def:physics:phyx-mini-0912:phyxminiproblems-problemphyxmini0912-answerchargeincoulombs}
  A choice is at least as close to the physical charge as every alternative
\end{definition}
\begin{proof}
  This is the declared model definition of Is Closest Displayed Charge.
\end{proof}
\begin{lemma}[ball charge choice C rounding bounds]
  \label{lem:physics:phyx-mini-0912:phyxminiproblems-problemphyxmini0912-ball-charge-choicec-rounding-bounds}
  \lean{PhyXMiniProblems.ProblemPhyXMini0912.ball_charge_choiceC_rounding_bounds}
  \uses{def:physics:phyx-mini-0912:phyxminiproblems-problemphyxmini0912-chargeincoulombs, def:physics:phyx-mini-0912:phyxminiproblems-problemphyxmini0912-protonorbitsetup, def:physics:phyx-mini-0912:phyxminiproblems-problemphyxmini0912-matchesproblemreadouts, def:physics:phyx-mini-0912:phyxminiproblems-problemphyxmini0912-matchesprimaryorbitfigure, def:physics:phyx-mini-0912:phyxminiproblems-problemphyxmini0912-satisfiesorbitradiusgeometry, def:physics:phyx-mini-0912:phyxminiproblems-problemphyxmini0912-usestextbookprotonandcoulombconstants, def:physics:phyx-mini-0912:phyxminiproblems-problemphyxmini0912-hasphysicalprotonorbitparameters, def:physics:phyx-mini-0912:phyxminiproblems-problemphyxmini0912-satisfiesexteriorsphericalcoulombfieldlaw, def:physics:phyx-mini-0912:phyxminiproblems-problemphyxmini0912-satisfieselectricforcelawonproton, def:physics:phyx-mini-0912:phyxminiproblems-problemphyxmini0912-satisfiesuniformcircularorbitlaws}
  The textbook constants and orbit data put the charge between -9.95 * 10\textasciicircum{}-12 C and -9.85 * 10\textasciicircum{}-12 C, so it displays as -9.9 * 10\textasciicircum{}-12 C
\end{lemma}
\begin{proof}
  The conclusion follows from the governing relations and figure readouts named in the statement of ball charge choice C rounding bounds.
\end{proof}
% --- Archon declaration topology end ---
% --- Archon physics formalization source end ---
